\section{本章重點}
原子作為自然界組成的基礎，在軌域中，需要了解決定單質子原子(氫)以及多質子原子電子分佈型態的要素，諸如能量、自旋等等，以一言以蔽之，本章要回答的問題是：原子究竟呈現什麼樣的結構與形狀。

選修物理第五冊提到波耳提出氫原子模型，該模型讓我們能預測單電子原子的電子分佈。但在量子力學的出現，讓我們能預測多電子原子的分佈(s、p、d、f....)。